% GENERATED FILE. Edit canonical lessons, then run npm run generate:books.
% canonical-source-hash: fnv1a64:694080e5a83c4c5f
% canonical-lessons: PT-C296-algemas, PT-C296-algemar, PT-C296-prender, PT-C296-preso, PT-C296-cela

\chapter{Algemas e celas}
\label{ch:pt-algemas-e-celas}
\hlchaptermodality{296}

\begin{chapteropening}
\textbf{By the end of this chapter:} \emph{I can describe an arrest and life in a cell: handcuffs, prisoners and cells.}

Complete the last lesson of chapter 296: I can describe an arrest and life in a cell: handcuffs, prisoners and cells.
\end{chapteropening}

\section[as algemas]{as algemas — the handcuffs}

\label{lesson:PT-C296-algemas}



\begin{quote}
Before the new one: say the Portuguese for the arrest, then the Portuguese for the person arrested.
\end{quote}

\subsection*{You'll want to know: as algemas}
\textbf{as algemas} — \textquotedblleft{}the handcuffs\textquotedblright{}.

\begin{quote}
O agente pôs as algemas ao suspeito.
\end{quote}

\begin{quote}
\emph{The officer put the handcuffs on the suspect.}
\end{quote}

Usually plural; one cuff is \emph{a algema}.

\begin{cousinweb}
\emph{algema} comes from Arabic \emph{al-jāmiʿa} \textquotedblleft{}the fetter\textquotedblright{}, literally \textquotedblleft{}that which joins\textquotedblright{} (EtymDB).
\end{cousinweb}

\begin{grammarlens}[title={What you've built}]
One more word for telling, explaining and describing at B1.
\end{grammarlens}

\subsection*{Your turn}
\begin{itemize}
\raggedright
  \item \emph{Say it:} \emph{as algemas}
  \item \emph{Say it:} \emph{as algemas}, once more
  \item \emph{Say it:} a sentence of your own with \emph{as algemas}
  \item \emph{Recall:} say the Portuguese for the arrest, then the Portuguese for the person arrested, then say \emph{as algemas} again
\end{itemize}

\begin{tcolorbox}[breakable,colback=teal!4,colframe=teal!35!black,title={Before you move on}]
What does as algemas mean? (\textquotedblleft{}The handcuffs\textquotedblright{}.) Say it once more.
\end{tcolorbox}
\section[algemar]{algemar — to handcuff}

\label{lesson:PT-C296-algemar}



\begin{quote}
Before the new one: say the Portuguese for the person arrested, then the Portuguese for the handcuffs.
\end{quote}

\subsection*{You'll want to know: algemar}
\textbf{algemar} — \textquotedblleft{}to handcuff\textquotedblright{}.

\begin{quote}
Algemaram o homem dentro do carro.
\end{quote}

\begin{quote}
\emph{They handcuffed the man inside the car.}
\end{quote}

Built on \emph{as algemas}.

\begin{cousinweb}
\emph{algemar} is built on \emph{algema} \textquotedblleft{}a handcuff\textquotedblright{}, which comes from Arabic \emph{al-jāmiʿa} \textquotedblleft{}the fetter\textquotedblright{} (EtymDB, under \emph{algema}).
\end{cousinweb}

\begin{grammarlens}[title={What you've built}]
One more word for telling, explaining and describing at B1.
\end{grammarlens}

\subsection*{Your turn}
\begin{itemize}
\raggedright
  \item \emph{Say it:} \emph{algemar}
  \item \emph{Say it:} \emph{algemar}, once more
  \item \emph{Say it:} a sentence of your own with \emph{algemar}
  \item \emph{Recall:} say the Portuguese for the person arrested, then the Portuguese for the handcuffs, then say \emph{algemar} again
\end{itemize}

\begin{tcolorbox}[breakable,colback=teal!4,colframe=teal!35!black,title={Before you move on}]
What does algemar mean? (\textquotedblleft{}To handcuff\textquotedblright{}.) Say it once more.
\end{tcolorbox}
\section[prender]{prender — to catch, to arrest; to fasten}

\label{lesson:PT-C296-prender}



\begin{quote}
Before the new one: say the Portuguese for the handcuffs, then the Portuguese for to handcuff.
\end{quote}

\subsection*{You'll want to know: prender}
\textbf{prender} — \textquotedblleft{}to catch, to arrest; to fasten\textquotedblright{}.

\begin{quote}
A polícia prendeu o assaltante.
\end{quote}

\begin{quote}
\emph{The police caught the robber.}
\end{quote}

Also ‘to tie up, to fasten’: \emph{prender o cabelo}.

\begin{cousinweb}
\emph{prender} is inherited from Latin \emph{prehendere} \textquotedblleft{}to grasp, to seize\textquotedblright{} (EtymDB).
\end{cousinweb}

\begin{grammarlens}[title={What you've built}]
One more word for telling, explaining and describing at B1.
\end{grammarlens}

\subsection*{Your turn}
\begin{itemize}
\raggedright
  \item \emph{Say it:} \emph{prender}
  \item \emph{Say it:} \emph{prender}, once more
  \item \emph{Say it:} a sentence of your own with \emph{prender}
  \item \emph{Recall:} say the Portuguese for the handcuffs, then the Portuguese for to handcuff, then say \emph{prender} again
\end{itemize}

\begin{tcolorbox}[breakable,colback=teal!4,colframe=teal!35!black,title={Before you move on}]
What does prender mean? (\textquotedblleft{}To catch, to arrest; to fasten\textquotedblright{}.) Say it once more.
\end{tcolorbox}
\section[o preso]{o preso — the prisoner}

\label{lesson:PT-C296-preso}



\begin{quote}
Before the new one: say the Portuguese for to handcuff, then the Portuguese for to catch.
\end{quote}

\subsection*{You'll want to know: o preso}
\textbf{o preso} — \textquotedblleft{}the prisoner\textquotedblright{}.

\begin{quote}
O preso escreve cartas à mãe.
\end{quote}

\begin{quote}
\emph{The prisoner writes letters to his mother.}
\end{quote}

As an adjective: \emph{estar preso} — to be in jail, or to be stuck.

\begin{cousinweb}
\emph{preso} comes from Latin \emph{prehēnsus} \textquotedblleft{}seized\textquotedblright{} (EtymDB), the past participle of \emph{prehendere}.
\end{cousinweb}

\begin{grammarlens}[title={What you've built}]
One more word for telling, explaining and describing at B1.
\end{grammarlens}

\subsection*{Your turn}
\begin{itemize}
\raggedright
  \item \emph{Say it:} \emph{o preso}
  \item \emph{Say it:} \emph{o preso}, once more
  \item \emph{Say it:} a sentence of your own with \emph{o preso}
  \item \emph{Recall:} say the Portuguese for to handcuff, then the Portuguese for to catch, then say \emph{o preso} again
\end{itemize}

\begin{tcolorbox}[breakable,colback=teal!4,colframe=teal!35!black,title={Before you move on}]
What does o preso mean? (\textquotedblleft{}The prisoner\textquotedblright{}.) Say it once more.
\end{tcolorbox}
\section[a cela]{a cela — the cell}

\label{lesson:PT-C296-cela}



\begin{quote}
Before the new one: say the Portuguese for to catch, then the Portuguese for the prisoner.
\end{quote}

\subsection*{You'll want to know: a cela}
\textbf{a cela} — \textquotedblleft{}the cell\textquotedblright{}.

\begin{quote}
A cela era pequena e fria.
\end{quote}

\begin{quote}
\emph{The cell was small and cold.}
\end{quote}

Also a monk’s cell.

\begin{cousinweb}
\emph{cela} is inherited from Latin \emph{cella} \textquotedblleft{}a small room, a storeroom\textquotedblright{} (EtymDB).
\end{cousinweb}

\begin{grammarlens}[title={What you've built}]
One more word for telling, explaining and describing at B1.
\end{grammarlens}

\subsection*{Your turn}
\begin{itemize}
\raggedright
  \item \emph{Say it:} \emph{a cela}
  \item \emph{Say it:} \emph{a cela}, once more
  \item \emph{Say it:} a sentence of your own with \emph{a cela}
  \item \emph{Recall:} say the Portuguese for to catch, then the Portuguese for the prisoner, then say \emph{a cela} again
\end{itemize}

\begin{tcolorbox}[breakable,colback=teal!4,colframe=teal!35!black,title={Before you move on}]
What does a cela mean? (\textquotedblleft{}The cell\textquotedblright{}.) Say it once more.
\end{tcolorbox}
